%----------------------------------------------------------------------------------------
% NOTATIE EN FIGUREN VAN DE CURSUS FUNCTIES (2de graad)
% Staat in de preambule en niet in een hoofdstuk: ook de lijsten met oplossingen gebruiken dit.
% Vertrokken van notatie.tex van Integralen; de intervallen en de getallenas zoals in Eerstegraadsvergelijkingen.
%----------------------------------------------------------------------------------------
% Bewaarde grafiek, geplaatst door maak-pdf.py: bestand en pagina in figuren-cache/
\newcommand{\grafiekbestand}[2]{\includegraphics[page=#2,trim=3bp 3bp 3bp 3bp]{figuren-cache/#1}}

\providecommand{\R}{\mathbb{R}}
\newcommand{\dom}{\operatorname{dom}}
\newcommand{\ber}{\operatorname{ber}}
% Absolute waarde: \abs{-3}
\providecommand{\abs}[1]{\lvert #1\rvert}

% Intervallen (Vlaamse notatie): \ivgg{a}{b} = [a, b], \ivoo{a}{b} = ]a, b[, \ivgo{a}{b} = [a, b[, \ivog{a}{b} = ]a, b]
\newcommand{\ivgg}[2]{\mathopen{[}{#1}{,}\;{#2}\mathclose{]}}
\newcommand{\ivoo}[2]{\mathopen{]}{#1}{,}\;{#2}\mathclose{[}}
\newcommand{\ivgo}[2]{\mathopen{[}{#1}{,}\;{#2}\mathclose{[}}
\newcommand{\ivog}[2]{\mathopen{]}{#1}{,}\;{#2}\mathclose{]}}
\newcommand{\pinf}{{+\infty}}
\newcommand{\minf}{{-\infty}}

% Zoals stappen, maar met een bredere kolom voor de berekening: voor lange formules
\newenvironment{stappenbreed}{%
  \par\renewcommand{\arraystretch}{1.25}%
  \noindent\begin{tabular}{@{}>{$\displaystyle}p{9.8cm}<{$}@{\hspace{4mm}}>{\RaggedRight\small\itshape\arraybackslash}p{4.9cm}@{}}%
}{\end{tabular}\par}

% Zoals stappen, met eigen kolombreedtes: \begin{stappenvrij}{6.2cm}{3.8cm}. Voor naast een figuur.
\newenvironment{stappenvrij}[2]{%
  \par\renewcommand{\arraystretch}{1.25}%
  \noindent\begin{tabular}{@{}>{$\displaystyle}p{#1}<{$}@{\hspace{3mm}}>{\RaggedRight\small\itshape\arraybackslash}p{#2}@{}}%
}{\end{tabular}\par}

% Tekst of berekening links, figuur rechts: \naast{links}{rechts}
\newsavebox{\naastbox}
\newcommand{\naast}[3][]{%
  \par\noindent\sbox{\naastbox}{#3}%
  \begin{minipage}[t]{\dimexpr\linewidth-\wd\naastbox-3mm\relax}\vspace{0pt}#2\end{minipage}\hfill
  \begin{minipage}[t]{\wd\naastbox}\vspace{0pt}\usebox{\naastbox}\end{minipage}\par\addvspace{1mm}}
% Een rij figuren, gecentreerd
\newcommand{\figuren}[1]{\par\addvspace{1mm}{\centering #1\par}\addvspace{1.5mm}}
% Opgave met een figuur rechts: \begin{metgrafiek}{tikz} opgave en deelvragen \end{metgrafiek}
\newsavebox{\grafiekbox}
\newenvironment{metgrafiek}[2][]{%
  \sbox{\grafiekbox}{#2}%
  \begin{minipage}[t]{\dimexpr\linewidth-\wd\grafiekbox-3mm\relax}%
}{%
  \end{minipage}\hfill
  \begin{minipage}[t]{\wd\grafiekbox}\vspace{-2mm}\usebox{\grafiekbox}\end{minipage}\par}

%----------------------------------------------------------------------------------------
% TABELLEN
%----------------------------------------------------------------------------------------
% Tabel met waarden, horizontaal en met een raster: \waardetabel{aantal kolommen}{x-waarden}{y-waarden},
% de waarden gescheiden door &, in wiskundemodus. Met eigen namen: \waardetabel[$t$ (min)][$T$ (°C)]{5}{...}{...}
\NewDocumentCommand{\waardetabel}{O{$x$} O{$f(x)$} m m m}{%
  {\renewcommand{\arraystretch}{1.35}%
  \begin{tabular}{|>{\columncolor{accent!12}}c|*{#3}{>{$}c<{$}|}}\hline #1 & #4 \\ \hline #2 & #5 \\ \hline\end{tabular}}}
% Tekenschema: \tekenschema{aantal kolommen}{x-rij}{tekenrij}, bv. \tekenschema{3}{ & 2 & }{- & 0 & +}
% Met een andere naam dan f(x): \tekenschema[g(x)]{3}{...}{...}
\NewDocumentCommand{\tekenschema}{O{f(x)} m m m}{%
  {\renewcommand{\arraystretch}{1.35}%
  \begin{tabular}{>{$}c<{$}|*{#2}{>{$}c<{$}}}x & #3 \\ \hline #1 & #4\end{tabular}}}
% Verlooptabel (stijgen en dalen): zoals \tekenschema, met \stijgt, \daalt en \cst in de onderste rij
\newcommand{\stijgt}{\nearrow}
\newcommand{\daalt}{\searrow}
\newcommand{\cst}{\rightarrow}
% Een cel met max of min onder een functiewaarde: \extr{4}{max}
\newcommand{\extr}[2]{\underset{\text{\scriptsize #2}}{#1}}

%----------------------------------------------------------------------------------------
% GRAFIEKEN, binnen een tikzpicture. Zet de schaal met x=... en y=... (bv. [x=.5cm,y=.5cm]).
%----------------------------------------------------------------------------------------
% Rooster met assen: \rooster{xmin}{xstap}{xmax}{ymin}{ystap}{ymax}, met getallen bij elke lijn (niet bij 0).
% Andere namen voor de assen: \rooster[$t$][$h$]{...}. \roosterleeg doet hetzelfde zonder getallen.
% Getallen bij een grotere stap: \roosterleeg, en dan \xgetallen{2,4,6} en \ygetallen{10,20}.
% Getallen met een eigen label: \xlabels{0.5/{0,5},1.5/{1,5}} en \ylabels{...}.
\NewDocumentCommand{\roosterleeg}{O{$x$} O{$y$} m m m m m m}{%
  \draw[black!14,very thin,xstep=#4,ystep=#7] (#3,#6) grid (#5,#8);
  \draw[->,grijs] (#3,0) -- ({#5+.45*#4},0) node[right,text=black,font=\footnotesize] {#1};
  \draw[->,grijs] (0,#6) -- (0,{#8+.45*#7}) node[above,text=black,font=\footnotesize] {#2};}
\NewDocumentCommand{\rooster}{O{$x$} O{$y$} m m m m m m}{%
  \roosterleeg[#1][#2]{#3}{#4}{#5}{#6}{#7}{#8}
  \pgfmathsetmacro{\rtweex}{#3+#4}\pgfmathsetmacro{\rtweey}{#6+#7}
  \foreach \t in {#3,\rtweex,...,#5}{\pgfmathtruncatemacro{\tt}{round(\t)}\ifdim\t pt=0pt\else\node[below,inner sep=1.5pt,text=black,font=\scriptsize] at (\t,0) {$\tt$};\fi}
  \foreach \t in {#6,\rtweey,...,#8}{\pgfmathtruncatemacro{\tt}{round(\t)}\ifdim\t pt=0pt\else\node[left,inner sep=1.5pt,text=black,font=\scriptsize] at (0,\t) {$\tt$};\fi}}
\newcommand{\xgetallen}[1]{\foreach \t in {#1}{\node[below,inner sep=1.5pt,text=black,font=\scriptsize] at (\t,0) {$\t$};}}
\newcommand{\ygetallen}[1]{\foreach \t in {#1}{\node[left,inner sep=1.5pt,text=black,font=\scriptsize] at (0,\t) {$\t$};}}
\newcommand{\xlabels}[1]{\foreach \t/\l in {#1}{\node[below,inner sep=1.5pt,text=black,font=\scriptsize] at (\t,0) {$\l$};}}
\newcommand{\ylabels}[1]{\foreach \t/\l in {#1}{\node[left,inner sep=1.5pt,text=black,font=\scriptsize] at (0,\t) {$\l$};}}
% De oorsprong: \oorsprong zet een 0 linksonder de oorsprong
\newcommand{\oorsprong}{\node[below left,inner sep=1.5pt,text=black,font=\scriptsize] at (0,0) {$0$};}

% Een grafiek: \kromme[opties]{xmin}{xmax}{voorschrift in \x}, bv. \kromme{-2}{3}{(\x)^2-1}
\newcommand{\kromme}[4][]{\draw[accent,thick,#1] plot[domain=#2:#3,samples=90,smooth] (\x,{#4});}
% Een rechte y = mx + q: \rechte[opties]{xmin}{xmax}{m}{q}
\newcommand{\rechte}[5][]{\draw[accent,thick,#1] (#2,{(#4)*(#2)+(#5)}) -- (#3,{(#4)*(#3)+(#5)});}
% Een grafiek door gegeven punten, vloeiend: \doorpunten[opties]{(0,1) (1,3) (2,2)}
\newcommand{\doorpunten}[2][]{\draw[accent,thick,#1] plot[smooth,tension=.6] coordinates {#2};}
% Een gebroken lijn door gegeven punten: \gebroken[opties]{(0,1) (1,3) (2,2)}
\newcommand{\gebroken}[2][]{\draw[accent,thick,line join=round,#1] plot coordinates {#2};}
% Een punt: \punt{x}{y}, een open bolletje (hoort er niet bij): \openpunt{x}{y}; kleur met de optie
\newcommand{\punt}[3][accent]{\fill[#1] ({#2},{#3}) circle (2.2pt);}
\newcommand{\openpunt}[3][accent]{\filldraw[#1,fill=white,thick] ({#2},{#3}) circle (2.2pt);}
% Stippellijnen om een punt af te lezen: \aflezen{x}{y} (van de x-as naar het punt en naar de y-as)
\newcommand{\aflezen}[3][grijs]{\draw[#1,dashed,thin] ({#2},0) -- ({#2},{#3}) -- (0,{#3});}
% Het trapje van de richtingscoëfficiënt: \trapje{x0}{y0}{dx}{dy}{label dx}{label dy}
% van (x0, y0) eerst dx opzij, dan dy omhoog (of omlaag als dy negatief is)
\newcommand{\trapje}[6]{%
  \draw[inkt,thick,dashed,->,>=stealth] (#1,#2) -- ({#1+#3},#2);
  \draw[inkt,thick,dashed,->,>=stealth] ({#1+#3},#2) -- ({#1+#3},{#2+#4});
  \node[font=\scriptsize,text=inkt,anchor=\ifdim#4pt<0pt south\else north\fi,inner sep=2pt] at ({#1+#3/2},#2) {#5};
  \node[font=\scriptsize,text=inkt,anchor=\ifdim#3pt<0pt east\else west\fi,inner sep=2pt] at ({#1+#3},{#2+#4/2}) {#6};}
% Een tweede grafiek in dezelfde figuur: stijl tweede (gestreept, donker), en hulplijn voor y = c
\tikzset{tweede/.style={inkt,thick,dashed},hulplijn/.style={grijs,thick,densely dotted}}
% Een deel van de x-as in kleur (de oplossing van een ongelijkheid): \asdeel{a}{b}
\newcommand{\asdeel}[2]{\draw[accent,line width=2.6pt,opacity=.55] (#1,0) -- (#2,0);}

% GETALLENAS met haken voor intervallen (zoals in Eerstegraadsvergelijkingen):
% \getallenas{min}{max}{lijst} met de lijst als t/label, bv. {-2/-2,0/0,3/3}
% Een interval erop: \ivdeel{a}{b} en de haken \haakl{a} ([) en \haakr{b} (]).
% Een gesloten grens krijgt een haak naar het interval toe, een open grens een haak ervan weg.
\newcommand{\getallenas}[3]{%
  \draw[->,grijs] ({#1-.4},0) -- ({#2+.5},0);
  \foreach \t/\l in {#3}{\draw[grijs] (\t,.1) -- (\t,-.1) node[below=1pt,text=black,font=\scriptsize] {$\l$};}}
\newcommand{\ivdeel}[2]{\draw[accent,line width=2.4pt] (#1,0) -- (#2,0);}
\newcommand{\haakl}[1]{\draw[accent,line width=1.3pt] ({#1+.16},.26) -- (#1,.26) -- (#1,-.26) -- ({#1+.16},-.26);}
\newcommand{\haakr}[1]{\draw[accent,line width=1.3pt] ({#1-.16},.26) -- (#1,.26) -- (#1,-.26) -- ({#1-.16},-.26);}
